% =====================================================================
% structure-example.tex  —  shows the document shape + ONE full topic
% for studynotegeneratorcaveman. This is a PATTERN to copy, not a fixed
% template. Pair it with assets/preamble.tex.
%
% Full file would be:
%   \input{preamble.tex}
%   \begin{document}
%     <title block>          (below)
%     <intro "how to read" box>
%     \tableofcontents
%     <topic sections>       (one shown below)
%     <closing checklist>
%   \end{document}
% =====================================================================


% ----------------------- TITLE BLOCK -----------------------
\begin{center}
  {\Huge\bfseries\color{ink} Course Name — CT 1}\\[4pt]
  {\Large\bfseries\color{accent} Chapter 1: Chapter Title}\\[8pt]
  {\small Theory + Bangla intuition + caveman notes + solved math}\\[2pt]
  {\small\itshape Built from the student's own handwritten lecture notes — same flow}\\[6pt]
  {\color{accent}\rule{0.7\textwidth}{1pt}}
\end{center}

% ----------------------- "HOW TO READ" INTRO BOX -----------------------
\vspace{4pt}
\begin{tcolorbox}[colback=softgrey, colframe=ink, boxrule=0.6pt, arc=3pt, left=10pt, right=10pt]
{\bfseries\color{ink} এই নোট কীভাবে পড়বে:} প্রতিটা টপিকে আগে \kw{theory} (ইংরেজিতে, যেভাবে
exam-এ লিখবে), তারপর \textcolor{banglafr}{\textbf{সবুজ box}}-এ বাংলায় আসল ব্যাপারটা, তারপর
\textcolor{cavefr}{\textbf{কমলা box}}-এ ছোট্ট caveman ভাষায় মনে রাখার লাইন। নীল box = মুখস্থ
করার সূত্র, হলুদ box = exam hint।
\end{tcolorbox}

\tableofcontents
\vspace{6pt}\hrule\vspace{10pt}


% ===================== ONE FULL TOPIC (the pattern) =====================
\section{Refractive Index}

% (1) THEORY — English, exam-ready, math in math mode
The \kw{refractive index} $n$ is the ratio that tells how much light slows in a medium:
\[
n=\frac{c}{v},\qquad v=\frac{c}{\sqrt{\varepsilon_r}},\qquad \boxed{\,n=\sqrt{\varepsilon_r}\,}.
\]
The wave's $E$-field drives the bound electrons, creating an oscillating polarization that
re-radiates and effectively slows the wave — this slowing is what we observe as refraction.

% (4) KEY EQUATION BOX — isolate the formula to memorize
\begin{formbox}
\[
n=\frac{c}{v}=\sqrt{\varepsilon_r}\qquad\text{(non-magnetic medium, } \mu_r=1)
\]
\end{formbox}

% (5) EXAM BOX — use where the source flagged "exam:"
\begin{exambox}
\textbf{Q: Explain the speed using $\varepsilon_r$.}\\
From $v=\dfrac{c}{\sqrt{\varepsilon_r}}$: if $\varepsilon_r\!\uparrow$ then $v\!\downarrow$
(slower light, bigger $n$); if $\varepsilon_r\!\downarrow$ then $v\!\uparrow$ (faster light).
\end{exambox}

% (2) BENGALI INTUITION BOX — the physical picture, plain Bengali
\begin{banglabox}
$n=\dfrac{c}{v}$ — আলো কোনো মাধ্যমে ঢুকলে কতটা ধীর হয় তার মাপ। $E$-ক্ষেত্র পরমাণুর ইলেকট্রনকে
দোলায়, polarization হয়, আর সেই কারণে আলো ধীর হয় — এটাই refraction। $\varepsilon_r$ বড় হলে
বেশি বাধা, আলো ধীরে চলে; ছোট হলে আলো জোরে চলে। সূত্র: $n=\sqrt{\varepsilon_r}$।
\end{banglabox}

% (3) CAVEMAN BOX — 1-3 line hook, arrows for causality
\begin{cavebox}
$n=c/v$. $E$ shake electron $\rightarrow$ polarize $\rightarrow$ light slow = refraction.
Big $\varepsilon_r$ = big obstacle = slow light = big $n$. Remember: $n=\sqrt{\varepsilon_r}$.
\end{cavebox}

% (7) WORKED EXAMPLE — full steps, reuse the source/textbook example number
\subsection{Solved Example X.Y}
\textbf{Given} $\varepsilon_r=2.25$:
\[
n=\sqrt{2.25}=1.5,\qquad v=\frac{c}{n}=\frac{3\times10^{8}}{1.5}=2\times10^{8}\ \mathrm{m/s}.
\]


% ===================== CLOSING CHECKLIST =====================
\section*{\color{ink}CT-1 Topic Checklist}
\addcontentsline{toc}{section}{CT-1 Topic Checklist}
\begin{tcolorbox}[colback=softgrey, colframe=ink, boxrule=0.7pt, arc=3pt]
\begin{multicols}{2}
\begin{enumerate}
  \item Topic one
  \item Topic two
  \item Topic three
\end{enumerate}
\end{multicols}
\textbf{Solved-example targets:} list the example numbers the student must be able to do.
\end{tcolorbox}
